% Invented deck exercising the M5 slice end to end: title metadata and
% \maketitle, section dividers, frames with titles (argument and
% \frametitle forms), verbatim code, math environments and tables carried
% as source and rows, unknown wrapper environments degrading to their body,
% and beamer idioms through the compat layer. Nothing here derives from any
% private document.
\documentclass[aspectratio=169]{beamer}

\usetheme{example}
\setbeamercovered{transparent}
\newcommand{\dim}[1]{\textcolor{quiet!50}{#1}}
\newenvironment{banner}[1]{}{}

\fonts{
  dir = "fonts",
  body = "Open Sans",
  mono = "Source Code Pro",
}

\title{A Certified Deck}
\subtitle{Slides from one IR}
\author{Pat Placeholder}
\institute{example.org}
\date{}

\begin{document}

\maketitle

\section{Engine}

\begin{frame}{Two backends}
    \begin{itemize}
        \item<1-> HTML is the deck; PDF is the handout.
        \item<2-> \alert{Frames are page boundaries}, never flattened.
    \end{itemize}
    \note{Speaker notes are hidden until notes land.}
\end{frame}

\begin{banner}{decoration}
    \begin{frame}[fragile]{Code}
        Definition:
\begin{verbatim}
def pages (doc : Doc) : Nat :=
    (layout doc).pages.size
\end{verbatim}
        \uncover<2>{\dim{Checked at build time.}}
    \end{frame}
\end{banner}

\begin{frame}
    \frametitle{Arithmetic}
    The first squares:
    \begin{align*}
        1 &= 1^2
        \\1 + 3 &= 2^2
    \end{align*}

    \begin{tabular}{ll}
        \textbf{Pages} & \textbf{Frames} \\
        \midrule
        one & per frame \\
        one & per divider \\
    \end{tabular}
\end{frame}

\begin{frame}[standout]
    Questions?
\end{frame}

\end{document}
